\chapter{Responses as values and bytes}
\chapterintro{Handlers should describe a result; one serializer should enforce wire rules. Centralizing response construction prevents every route from reimplementing HTTP.}

\begin{objectives}
\begin{itemize}
  \item Model status, fields, and body independently of socket I/O.
  \item Serialize legal status lines and fields without response splitting.
  \item Implement Content-Length, HEAD, no-content statuses, and connection fields correctly.
  \item Keep internal error detail out of public 500 responses.
\end{itemize}
\end{objectives}

\section{The response object}

A handler should return a \code{Response}, not write to the socket. The response
contains an integer status, a multi-value \code{Headers} collection, and a bytes
body. This separation gives middleware a chance to inspect or replace the result,
lets tests call handlers without networking, and gives one serializer authority
over HTTP framing.

Stage 4 isolates the concept in a small program:

\codefile{../code/stages/04_response_serializer/server.py}

The stage already computes content length, uses the standard status phrase, and
sends no body for HEAD. The final serializer adds date/server fields, persistent
connection policy, special no-content statuses, and header validation.

\section{Headers must be safe to serialize}

Response splitting becomes possible when untrusted text inserts carriage return
or line feed into a field value. A value such as
\code{"safe\\r\\nX-Injected: yes"} would create a second field if serialized
blindly. The final \code{Headers.add} requires a legal token name and rejects
control characters in values.

This validation occurs when the collection is built, not after bytes have been
joined. As a result, every response producer shares the same invariant: header
items can be represented as one field line each.

\section{One serialization path}

The server's response method performs these steps:

\begin{enumerate}
  \item resolve the standard reason phrase for the numeric status;
  \item copy fields so server-added metadata does not mutate the handler's object;
  \item set Date, Server, and Connection;
  \item determine whether the status permits content;
  \item set Content-Length to the representation-body length when permitted;
  \item suppress transmitted bytes for HEAD while preserving that length;
  \item encode the start line as ASCII and fields as ISO-8859-1;
  \item join one blank CRLF line and the permitted body, then call \code{sendall}.
\end{enumerate}

The serializer operates on bytes at its final boundary. It does not call
\code{json.dumps}; JSON is a response factory concern.

\section{HEAD and no-content statuses differ}

For HEAD, the server computes the response that GET would have produced and
suppresses content bytes. A recipient can learn the media type and expected
length without transferring the representation.

For 204 and 304, the status semantics prohibit message content. The serializer
suppresses the body and removes \code{Content-Length} in this implementation.
A handler that mistakenly attaches bytes cannot violate the wire rule.

\begin{pitfall}
Do not change a HEAD response's Content-Length to zero merely because zero body
bytes are transmitted. That field describes the selected representation a GET
would have carried, enabling useful metadata checks.
\end{pitfall}

\section{JSON and problem factories}

\code{Response.json(value)} uses compact separators, preserves Unicode, encodes
UTF-8, and sets \code{application/json; charset=utf-8}. A central factory means
all handlers use the same encoding and media type.

\code{Response.problem} creates an RFC 9457-style problem document with
\code{type}, \code{title}, \code{status}, \code{detail}, optional
\code{instance}, and optional field errors. It changes the media type to
\code{application/problem+json}. Clients can handle errors structurally instead
of scraping prose.

An expected \code{HTTPError} carries public-safe detail. An unexpected exception
is logged internally and converted to a generic 500 problem. Stack traces,
database paths, SQL statements, and secrets must not cross the response boundary.

\begin{wireview}
An invalid task might receive:
\begin{wiremessage}
HTTP/1.1 422 Unprocessable Content
Content-Type: application/problem+json; charset=utf-8
Content-Length: ...

{"type":"urn:task-api:problem:validation-error",
 "title":"Validation failed","status":422,
 "detail":"One or more values failed validation.",
 "errors":{"title":"Must be a non-empty string."}}
\end{wiremessage}
The line wrapping above is for print; the actual JSON is one compact byte sequence.
\end{wireview}

\section{Test meaning and wire form separately}

Handler tests should assert response status, fields, and decoded JSON. Serializer
tests should assert exact CRLF framing, length, HEAD suppression, 204 behavior,
and connection policy. Keeping these tests separate reveals which abstraction
failed.

\summarytitle

Response values decouple application decisions from network serialization.
Centralized header validation and serialization enforce length, HEAD, no-content,
and connection rules consistently. JSON and problem factories produce predictable
media types while unexpected failures remain private.

\reviewtitle
\begin{enumerate}
  \item Why should a handler not write directly to a connection?
  \item How can a newline in a field value become a vulnerability?
  \item How do HEAD and 204 differ in their treatment of Content-Length?
  \item Which error information belongs in logs but not a 500 response?
\end{enumerate}

\begin{exercise}[Serializer matrix]
Create expected start line, Content-Length presence, and transmitted body length
for GET/200, HEAD/200, GET/204, GET/304, and GET/500. Turn the matrix into unit
tests before changing serializer code.
\end{exercise}

\begin{exercise}[Injection attempt]
Construct a response field whose value contains CRLF. Verify that \code{Headers}
rejects it before serialization. Explain why percent-encoding that input would
not be a general replacement for validation.
\end{exercise}

